% Provenance (never printed): see sources/outline_sources.md (Year 2: tholeiitic, calc-alkaline
% and alkaline series; the TAS diagram; settings; isotopic signatures of mantle and crust).
% Classification from BGS RR 99-06 (after Streckeisen 1976, Le Bas et al. 1986, Le Maitre et
% al. 1989); series and data from Vermeesch and Pease (2021) and the GeoplotR package;
% curves and counts from figdata/bachelor-2/series.py; exercise and problem values in the
% Book 4 scratch calculators calc/ch09.py and calc/ch09pb.py.
\chapter{岩浆系列与构造背景}\label{ch:b2:magmatic-series}

只给岩石学家一份化学分析而不给地图，她往往能判断
熔岩是在洋脊喷发、在洋岛喷发，还是在俯冲板块之上喷发。
冰岛的熔岩与北美洲西部喀斯喀特火山的熔岩，二者都
由玄武岩演化至流纹岩；成百上千地分析后发现，它们在演化时沿着两条不同的路径，
一条越来越富铁，另一条则不然。本章首先
为火成岩命名：矿物可辨时依据矿物，
矿物不可辨时依化学分析；随后把岩石归并为若干系列，即岩石家族，
这些家族由分异作用相互联系，每个系列对应一种构造背景，
并介绍岩石学家据分析判断构造背景所用的图解，以及这些图解的局限。
% ledger: vp21:series

\begin{recall}
\textbf{来自高中分册}：玄武岩、安山岩、流纹岩、辉长岩和花岗岩；
长英质岩与镁铁质岩；火山岩结构与深成岩结构。\textbf{来自本科一年级
分册}：洋中脊、俯冲带与火山弧、热点与大
火成岩省、大陆裂谷。\textbf{来自\cref{ch:b2:elements-in-the-earth}}：
MORB 与 OIB、\omterm{def:b2:elements-in-the-earth:bse}{原始地幔}标准化（蛛网）图、\omterm{def:b2:elements-in-the-earth:partition}{相容元素}和
\omterm{def:b2:elements-in-the-earth:partition}{不相容元素}。\textbf{来自\cref{ch:b2:isotope-geochemistry}}：
\omterm{def:b2:isotope-geochemistry:reservoirs}{亏损地幔}和富集同位素储库。\textbf{来自
\cref{ch:b2:magmatic-differentiation}}：\omterm{def:b2:magmatic-differentiation:crystallisation}{分离结晶}、\omterm{def:b2:magmatic-differentiation:lld}{液相演化线}、
\omterm{def:b2:magmatic-differentiation:harker}{哈克图解}与瑞利定律。
\end{recall}

\section{火成岩命名}\label{sec:b2:magmatic-series:naming}

粗粒岩石可见其矿物，据矿物命名。熔岩往往
过于细粒或玻璃质而无法如此命名，只能据化学分析命名。两种
方案都是国际通行的约定，它们只作描述，不暗示
岩石的形成过程。
% ledger: bgs:tas-use

\begin{definition}[实际矿物组成；QAPF 图解]\label{def:b2:magmatic-series:qapf}
岩石的\emph{实际矿物组成}\index{实际矿物组成}是指其中各矿物
所占体积的比例。\emph{QAPF 图解}\index{QAPF 图解}用于命名
镁铁质矿物含量低于\qty{90}{\percent}的粗粒火成岩，
由四种模式组分（重新折算为总和为 100）：石英（Q）、碱性长石
（A）、斜长石（P）和似长石类（F）。石英与似长石类不会
同时结晶，因此每种岩石落在两个三角形之一，即 QAP 或 FAP，二者
沿公共边 AP 相连。
\end{definition}
% ledger: bgs:qapf

\begin{omfigure}
\centering
\begin{tikzpicture}[font=\scriptsize, x=1cm, y=1cm]
\def\L{6.4}
\pgfmathsetmacro{\H}{\L*0.8660254}
% upper point (q, p): q = Q fraction, p = P/(A+P)
\newcommand{\qapfup}[2]{({(#1)*\L/2 + (#2)*(1-(#1))*\L}, {(#1)*\H})}
\newcommand{\dn}[2]{({(#1)*\L/2 + (#2)*(1-(#1))*\L}, {-(#1)*\H})}
\draw[thick] (0,0) -- (\L/2,\H) -- (\L,0) -- (\L/2,-\H) -- cycle;
\draw (0,0) -- (\L,0);
\foreach \q in {0.9,0.6,0.2,0.05} \draw \qapfup{\q}{0} -- \qapfup{\q}{1};
\foreach \q in {0.1,0.6} \draw \dn{\q}{0} -- \dn{\q}{1};
\foreach \p in {0.1,0.65,0.9} \draw \qapfup{0.6}{\p} -- \qapfup{0}{\p};
\draw[dashed] \qapfup{0.6}{0.35} -- \qapfup{0.2}{0.35};
\draw \qapfup{0.2}{0.35} -- \qapfup{0}{0.35};
\foreach \p in {0.1,0.35,0.65,0.9} \draw \dn{0}{\p} -- \dn{0.1}{\p};
\foreach \p in {0.1,0.5,0.9} \draw \dn{0.1}{\p} -- \dn{0.6}{\p};
\node[anchor=south] at (\L/2,\H) {\normalsize Q};
\node[anchor=east] at (0,0) {\normalsize A};
\node[anchor=west] at (\L,0) {\normalsize P};
\node[anchor=north] at (\L/2,-\H) {\normalsize F};
\node at \qapfup{0.94}{0.5} {1a};
\node at \qapfup{0.75}{0.5} {1b};
\node[font=\tiny] at \qapfup{0.4}{0.05} {2};
\node at \qapfup{0.4}{0.225} {3a};
\node at \qapfup{0.4}{0.5} {3b};
\node at \qapfup{0.4}{0.775} {4};
\node[font=\tiny] at \qapfup{0.4}{0.95} {5};
\foreach \p/\n in {0.225/7*,0.5/8*,0.775/9*} \node at \qapfup{0.125}{\p} {\n};
\node[font=\tiny] at \qapfup{0.125}{0.05} {6*};
\node[font=\tiny] at \qapfup{0.125}{0.95} {10*};
\foreach \p/\n in {0.225/7,0.5/8,0.775/9} \node at \qapfup{0.025}{\p} {\n};
\foreach \p/\n in {0.225/7',0.5/8',0.775/9'} \node at \dn{0.05}{\p} {\n};
\node[font=\tiny] at \dn{0.3}{0.05} {11};
\node at \dn{0.35}{0.3} {12};
\node at \dn{0.35}{0.7} {13};
\node[font=\tiny] at \dn{0.3}{0.95} {14};
\node at \dn{0.75}{0.5} {15};
\node[anchor=east] at \qapfup{0.9}{0} {90};
\node[anchor=east] at \qapfup{0.6}{0} {60};
\node[anchor=east] at \qapfup{0.2}{0} {20};
\node[anchor=east] at \dn{0.1}{0} {10};
\node[anchor=east] at \dn{0.6}{0} {60};
\foreach \p/\n in {0.1/10,0.35/35,0.65/65,0.9/90} \node[anchor=south, fill=white, inner sep=0.5pt] at \qapfup{0.6}{\p} {\tiny\n};
\node[anchor=west, align=left, font=\scriptsize] at (\L+0.9,1.6) {上三角形：\% Q\\在 Q + A + P 中；\\
  等值线：\\$100\,P/(A+P)$\\取 10、35、65、90};
\node[anchor=west, align=left, font=\scriptsize] at (\L+0.9,-1.7) {下三角形：\% F\\在 F + A + P 中；\\
  $100\,P/(A+P)$\\取 10、50、90};
\end{tikzpicture}
\par\smallskip{\small 粗粒火成岩的 QAPF 分类。1a
石英岩，1b 富石英花岗类岩，2 碱性长石花岗岩，3a 正长花岗岩和
3b 二长花岗岩（合称花岗岩），4 花岗闪长岩，5 英云闪长岩；6*--10* 为
6--10 的含石英变种：6 碱性长石正长岩，7 正长岩，8 二长岩，9 二长闪长岩
或二长辉长岩，10 闪长岩、辉长岩或斜长岩；6'--10' 的
含似长石类变种；11 含似长石类正长岩，12 含似长石类二长正长岩，13 含似长石类二长闪长岩或
二长辉长岩，14 含似长石类闪长岩或辉长岩，15 似长石岩。细窄区 6、6'、10 和 10'
沿边缘，未编号。分类依英国地质调查局方案（
基于 Streckeisen 1976）。}
\end{omfigure}
% ledger: bgs:qapf

10 区含三个名称。\omterm{def:b2:magmatic-series:qapf}{QAPF 图解}无法区分闪长岩、辉长岩和
斜长岩，也不能进一步区分辉长类岩石；另有依据
斜长石和镁铁质矿物的图解来完成此区分。

\begin{definition}[全碱–硅图解]\label{def:b2:magmatic-series:tas}
\emph{全碱–硅图解}\index{全碱–硅图解}（TAS
图解）依据火山岩的化学成分为其命名。这类岩石的矿物组成难以直接测定，故
以 \ce{Na2O + K2O} 对 \ce{SiO2} 的质量百分数为据，分析结果须先
换算为不含水和二氧化碳、合计为 \qty{100}{\percent} 的数据。图解以直线划分为十五个区，
这些直线的交点坐标固定，
（见下段之后的图）。
\end{definition}
% ledger: bgs:tas, bgs:tas-use

图解分为三组区域。底部为低碱区，列出常见系列的名称，即苦橄玄武岩、玄武岩至流纹岩，
以二氧化硅含量 45、52、
57 和 \qty{63}{\percent} 为界，流纹岩区另以斜线分隔。其上的 S 区与粗面岩容纳同一二氧化硅含量下碱含量更高的岩石，
U 区、响岩和似长石岩容纳的岩石碱含量极高，
二氧化硅极少，以致结晶出似长石类矿物。少数区
再以第二判据细分：粗面玄武岩若 $\ce{Na2O} - 2
\ge \ce{K2O}$ 则称夏威夷岩，否则为钾质粗面玄武岩；标准矿物中橄榄石含量超过
\qty{10}{\percent}（见下）的碧玄岩为碱玄岩；流纹岩或粗面岩，其
分子比 $(\ce{Na2O + K2O})/\ce{Al2O3}$ 大于 1 时称为过碱性。该方案
要求分析的含水量低于 \qty{2}{\percent}，
二氧化碳低于 \qty{0.5}{\percent}：蚀变会带入水并迁移碱质，\omterm{def:b2:magmatic-differentiation:cumulate}{堆晶岩}
并非液体。
% ledger: bgs:tas, bgs:tas-use

\begin{omfigure}
\centering
\begin{tikzpicture}
\begin{axis}[omrecord, width=0.86\linewidth, height=8.6cm,
  xmin=37, xmax=80, ymin=0, ymax=16, xtick={37,41,...,77}, ytick={0,2,...,16},
  xlabel={\ce{SiO2} (\unit{\percent})}, ylabel={\ce{Na2O + K2O} (\unit{\percent})},
  legend style={font=\scriptsize, at={(0.98,0.98)}, anchor=north east},
  tick label style={font=\scriptsize}, label style={font=\small}]
\addplot[only marks, mark=*, mark size=0.55pt, blue!70!black, opacity=0.6] table[x=SiO2, y=alk] {figdata/out/bachelor-2-series-cath-th.dat};
\addlegendentry{冰岛}
\addplot[only marks, mark=*, mark size=0.55pt, red!75!black, opacity=0.6] table[x=SiO2, y=alk] {figdata/out/bachelor-2-series-cath-ca.dat};
\addlegendentry{喀斯喀特山脉}
\foreach \f in {Pc,B,O1,O2,O3,R,S1,S2,S3,T,U1,U2,U3,Ph,F}
  \addplot[black, thin, forget plot] table[x=SiO2, y=alk] {figdata/out/bachelor-2-series-tas-\f.dat};
\node[font=\tiny] at (axis cs:43,1.6) {Pc};
\node[font=\tiny, fill=white, inner sep=0.5pt] at (axis cs:50.5,0.8) {B};
\node[font=\tiny] at (axis cs:54.5,2.2) {O1};
\node[font=\tiny] at (axis cs:60,2.2) {O2};
\node[font=\tiny] at (axis cs:66.5,2.2) {O3};
\node[font=\tiny] at (axis cs:74.5,11) {R};
\node[font=\tiny, fill=white, inner sep=0.5pt] at (axis cs:49,5.8) {S1};
\node[font=\tiny] at (axis cs:52.5,7.4) {S2};
\node[font=\tiny] at (axis cs:57.5,8.9) {S3};
\node[font=\tiny] at (axis cs:64.5,12) {T};
\node[font=\tiny] at (axis cs:44,6.6) {U1};
\node[font=\tiny] at (axis cs:48.7,9.4) {U2};
\node[font=\tiny] at (axis cs:52.5,11.8) {U3};
\node[font=\tiny] at (axis cs:57,14.6) {Ph};
\node[font=\tiny] at (axis cs:39,10) {F};
\end{axis}
\end{tikzpicture}
\par\smallskip{\small TAS 图解（据 Le Bas 等，1986，经
英国地质调查局给出）。冰岛（蓝色）的 456 个熔岩和侵入岩，
喀斯喀特山脉（红色）的 580 个，氧化物按分析值给出。Pc 苦橄玄武岩、B 玄武岩、O1 玄武安山岩、
O2 安山岩、O3 英安岩、R 流纹岩；S1 粗面玄武岩、S2 玄武质粗安岩、
S3 粗安岩，T 粗面岩或粗面英安岩；U1 碧玄岩或
碱玄岩，U2 响碧玄岩，U3 碧响岩，Ph 响岩，F 似长石岩。}
\end{omfigure}
% ledger: bgs:tas, geoplotr:tas, data:b4-cath


\begin{definition}[硅饱和]\label{def:b2:magmatic-series:saturation}
称\omterm{def:b2:magma-emplacement:magma}{岩浆}或岩石为\emph{硅过饱和}\index{硅饱和}，若其二氧化硅多于其他氧化物
能与硅酸盐矿物结合的量，以致出现石英（或其他二氧化硅矿物）；
若二氧化硅过少，则称为\emph{硅不饱和}，
即不足以使全部碱质转化为长石，以致出现似长石类矿物（霞石、
白榴石）代之，且绝不与石英同时出现；而\emph{硅饱和}
介于二者之间。
\end{definition}

对于矿物无法计数的熔岩，饱和性依据\emph{标准矿物}判断：
这是一种标准计算，将分析所得氧化物分配给理想矿物，
按固定顺序进行，先将碱质归入长石，多余的二氧化硅记为石英，
不足者记为似长石类矿物。常用的 CIPW 标准矿物计算法可区分两类玄武岩：
标准矿物中含霞石的玄武岩为碱性玄武岩；
不含霞石的为亚碱性玄武岩，其中拉斑玄武岩的标准矿物含
紫苏辉石（斜方辉石），实际矿物中常含低钙辉石。
\Cref{exo:b2:magmatic-series:11} 给出了它的简化版本。
% ledger: bgs:tas-use, bgs:series-terms

\section{岩浆系列}\label{sec:b2:magmatic-series:series}

同一火山或同一岩区的熔岩很少完全相同，但彼此相关：
它们源于相似的母\omterm{def:b2:magma-emplacement:magma}{岩浆}，经
\cref{ch:b2:magmatic-differentiation}中的分异作用形成，并共同追踪一条演化路线。

\begin{definition}[岩浆系列]\label{def:b2:magmatic-series:series}
称\emph{岩浆系列}\index{岩浆系列}为同一\omterm{def:b2:magma-emplacement:magma}{岩浆}旋回或同一岩区的一组火成岩，
由相似母\omterm{def:b2:magma-emplacement:magma}{岩浆}经分异作用相互关联，并
具有共同的化学和矿物学特征。通常可区分为三种。
其中，\emph{拉斑玄武岩系列}\index{拉斑玄武岩系列}始于
\omterm{def:b2:magmatic-series:saturation}{硅饱和}或略过饱和的玄武岩，标准矿物中含紫苏辉石，
且在分异之初显著富铁。
\emph{钙碱性系列}\index{钙碱性系列}由玄武岩经
安山岩演化至英安岩和流纹岩，富铁程度很低或无，并且
其演化程度较高的成员含角闪石、黑云母等含水矿物。\emph{碱性系列}\index{碱性系列}
始于标准矿物含霞石的玄武岩，
经 TAS 图解中的 S 区和 U 区，向粗面岩和响岩演化。
\end{definition}
% ledger: bgs:series-terms, bgs:tas-use

这些名称描述的是系列，而非单个岩石：孤立的一块玄武岩可能属于其中任一系列，
且同一名称曾被赋予不同含义，因此岩石学家
须说明其所指何义。\omterm{def:b2:magmatic-series:series}{拉斑玄武岩系列}与\omterm{def:b2:magmatic-series:series}{钙碱性系列}之间的区分
取决于铁，而最能显示这一点的图解
所投的是三组氧化物。
% ledger: bgs:series-terms

\begin{definition}[AFM 图解]\label{def:b2:magmatic-series:afm}
称 \emph{AFM 图解}\index{AFM 图解}为三角图解，其三个顶点为
$A = \ce{Na2O + K2O}$、$F = \ce{FeO_T}$（全部铁以 \ce{FeO} 计）和 $M = \ce{MgO}$，
各组分以质量百分数表示，并换算使 $A + F + M = 100$。图上一条线将
偏向 F 一侧的拉斑玄武岩区与钙碱性区分开。
\end{definition}
% ledger: vp21:series, geoplotr:afm

分异中的玄武岩流失镁最快，因为橄榄石和辉石中的镁比液体中的多，
而碱质被保留在\omterm{def:b2:magma-emplacement:magma}{岩浆}中，因为它们是不相容的：
其演化线远离 M，朝向 A。铁的去向
决定了系列。在拉斑质\omterm{def:b2:magma-emplacement:magma}{岩浆}中，铁只随富镁的
亚铁硅酸盐析出，且比镁慢，因此比值
$F/M$ 上升，演化线先朝 F 攀升，随后转向 A：
即\emph{芬纳趋势}。在钙碱性\omterm{def:b2:magma-emplacement:magma}{岩浆}中，其氧化程度更高，铁还随
磁铁矿析出，磁铁矿含三价铁，$F/M$ 在液体
径直朝 A 移动时几乎保持不变：即\emph{鲍文趋势}。下图以真实岩石展示两者：
冰岛的拉斑玄武岩与喀斯喀特的钙碱性熔岩。
分界线由 1971 年目测绘制，在岩石为玄武质时能很好地将二者分开；
碱含量高时，两条趋势趋于汇合，分界线几乎无所作为。
% ledger: vp21:series, geoplotr:afm, data:b4-cath

\begin{omfigure}
\centering
\begin{tikzpicture}
\begin{axis}[name=afm, width=0.5\linewidth, height=0.5\linewidth, axis equal image,
  hide axis, xmin=-0.04, xmax=1.04, ymin=-0.06, ymax=0.92, clip=false]
\addplot[only marks, mark=*, mark size=0.5pt, blue!70!black, opacity=0.6] table[x=afm_x, y=afm_y] {figdata/out/bachelor-2-series-cath-th.dat};
\addplot[only marks, mark=*, mark size=0.5pt, red!75!black, opacity=0.6] table[x=afm_x, y=afm_y] {figdata/out/bachelor-2-series-cath-ca.dat};
\addplot[very thick, black] table[x=x, y=y] {figdata/out/bachelor-2-series-afm-line.dat};
\draw[thick] (axis cs:0,0) -- (axis cs:1,0) -- (axis cs:0.5,0.8660254) -- cycle;
\node[anchor=south] at (axis cs:0.5,0.866) {F};
\node[anchor=north east] at (axis cs:0,0) {A};
\node[anchor=north west] at (axis cs:1,0) {M};
\node[font=\scriptsize, fill=white, inner sep=1pt] at (axis cs:0.52,0.62) {拉斑玄武岩};
\node[font=\scriptsize, fill=white, inner sep=1pt] at (axis cs:0.33,0.12) {钙碱性};
\end{axis}
\begin{axis}[omrecord, at={(afm.east)}, anchor=west, xshift=1.3cm,
  width=0.46\linewidth, height=0.46\linewidth,
  xmin=-2.5, xmax=2.2, ymin=-4.5, ymax=1.5,
  xlabel={$\ln(A/F)$}, ylabel={$\ln(M/F)$},
  legend style={font=\scriptsize, at={(0.98,0.98)}, anchor=north east},
  tick label style={font=\scriptsize}, label style={font=\small}]
\addplot[only marks, mark=*, mark size=0.5pt, blue!70!black, opacity=0.6] table[x=lnAF, y=lnMF] {figdata/out/bachelor-2-series-lr-th.dat};
\addlegendentry{冰岛}
\addplot[only marks, mark=*, mark size=0.5pt, red!75!black, opacity=0.6] table[x=lnAF, y=lnMF] {figdata/out/bachelor-2-series-lr-ca.dat};
\addlegendentry{喀斯喀特}
\addplot[thick, blue!50!black, forget plot] table[x=x, y=y_th] {figdata/out/bachelor-2-series-lr-fits.dat};
\addplot[thick, red!50!black, forget plot] table[x=x, y=y_ca] {figdata/out/bachelor-2-series-lr-fits.dat};
\end{axis}
\end{tikzpicture}
\par\smallskip{\small 左：TAS 图中冰岛（蓝色）与喀斯喀特（红色）
岩石的 \omterm{def:b2:magmatic-series:afm}{AFM 图解}，附 Irvine 与 Baragar（1971）的分界线。右：
同一批岩石的对数比坐标图，每个系列在其中近似为一条直带；
图中的线为正交回归，斜率分别为 $-1.10$（冰岛）和 $-0.72$
（喀斯喀特），两线相交于 $(-3.84, 2.12)$，与 Vermeesch 与 Pease（2021）所得一致。}
\end{omfigure}
% ledger: geoplotr:afm, data:b4-cath, vp21:model

在分界线范围内的冰岛岩石中，351 个落在拉斑一侧，
23 个落在另一侧；在喀斯喀特岩石中，417 个落在钙碱性一侧，
57 个落在另一侧。曲线作为边界很笨拙，简单的模型即可解释
趋势的形状。设\omterm{def:b2:magma-emplacement:magma}{岩浆}中的碱、铁和镁
以与其含量成正比的速率流失，即 $\mathrm dX/\mathrm dt = -\lambda_X X$。

\begin{proposition}[对数比趋势]\label{prop:b2:magmatic-series:logratio}
若 $A$、$F$ 和 $M$ 以常数 $\lambda_A$、$\lambda_F$ 和
$\lambda_M$ 指数衰减，则\omterm{def:b2:magma-emplacement:magma}{岩浆}在 $\ln(A/F)$ 与 $\ln(M/F)$ 平面中沿一条直线运动，
其斜率为 $(\lambda_F - \lambda_M)/(\lambda_F - \lambda_A)$，而在
AFM 三角形上则沿一条曲线运动。
\end{proposition}

\begin{proof}
由 $X = X_0 e^{-\lambda_X t}$，$\ln(A/F) = \ln(A_0/F_0) + (\lambda_F - \lambda_A)t$
且 $\ln(M/F) = \ln(M_0/F_0) + (\lambda_F - \lambda_M)t$：二者均为 $t$ 的线性函数，
故其一为另一的线性函数，斜率如上所述。比值 $A/F$ 与 $M/F$ 确定
三角形上的点，其坐标 $A/(A + F + M)$、\dots\ 对 $t$ 并非线性。
\end{proof}

冰岛岩石较陡的斜率意味着较小的 $\lambda_F$：铁离开
拉斑质\omterm{def:b2:magma-emplacement:magma}{岩浆}的速度较慢（见\cref{exo:b2:magmatic-series:10}）。Vermeesch 与 Pease 证明，上图中的两条
带相交于一点，他们将其解读为假想的共同
源区。
% ledger: vp21:model

拉斑系列与\omterm{def:b2:magmatic-series:series}{钙碱性系列}在钾含量上也有差异。
\ce{K2O}--\ce{SiO2} 图解将玄武岩、安山岩、英安岩和流纹岩
中二氧化硅含量超过 \qty{48}{\percent} 者，按两条直线
划分为低钾、中钾和高钾三类（见下图）。两个岩套的大部分在此均属中钾类：
该图说明一个系列含钾多少，而非它属于哪个系列。
% ledger: bgs:k2o-sio2, data:b4-cath

\begin{omfigure}
\centering
\begin{tikzpicture}
\begin{axis}[omrecord, width=0.78\linewidth, height=6.2cm,
  xmin=45, xmax=79, ymin=0, ymax=5,
  xlabel={\ce{SiO2} (\unit{\percent})}, ylabel={\ce{K2O} (\unit{\percent})},
  legend style={font=\scriptsize, at={(0.02,0.98)}, anchor=north west},
  tick label style={font=\scriptsize}, label style={font=\small}]
\addplot[only marks, mark=*, mark size=0.55pt, blue!70!black, opacity=0.6] table[x=SiO2, y=K2O] {figdata/out/bachelor-2-series-cath-th.dat};
\addlegendentry{冰岛}
\addplot[only marks, mark=*, mark size=0.55pt, red!75!black, opacity=0.6] table[x=SiO2, y=K2O] {figdata/out/bachelor-2-series-cath-ca.dat};
\addlegendentry{喀斯喀特}
\addplot[thick, black, forget plot] table[x=x, y=y_low] {figdata/out/bachelor-2-series-k2o-lines.dat};
\addplot[thick, black, forget plot] table[x=x, y=y_high] {figdata/out/bachelor-2-series-k2o-lines.dat};
\draw[thick] (axis cs:48,0) -- (axis cs:48,0.3);
\node[font=\scriptsize, fill=white, inner sep=1pt] at (axis cs:60,3.7) {高钾};
\node[font=\scriptsize, fill=white, inner sep=1pt] at (axis cs:76.5,2.2) {中钾};
\node[font=\scriptsize, fill=white, inner sep=1pt] at (axis cs:60,0.35) {低钾};
\end{axis}
\end{tikzpicture}
\par\smallskip{\small 冰岛与喀斯喀特岩石的钾含量对二氧化硅含量图，
附两条线，分别经过 (48, 0.3)--(68, 1.2) 与 (48, 1.2)--(68, 3.1)，用以划分低钾、中钾和高钾
类型（据 Peccerillo 与泰勒，1976，为国际方案所采用），
此处延伸至 68 以外。}
\end{omfigure}
% ledger: bgs:k2o-sio2, data:b4-cath

为何氧化程度不同？钙碱性\omterm{def:b2:magma-emplacement:magma}{岩浆}形成于俯冲带之上，
脱水的板片向地幔楔输送氧化态物种；\omterm{def:b2:magma-emplacement:magma}{岩浆}
结晶出磁铁矿，其演化较充分的成员则结晶出普通角闪石或黑云母。
拉斑质\omterm{def:b2:magma-emplacement:magma}{岩浆}形成于洋脊，那里地幔通过减压而熔融
（\cref{ch:b2:mantle-melting}），且其还原性更强。
% ledger: vp21:series

\section{岩浆系列与构造环境}\label{sec:b2:magmatic-series:settings}

每种板块构造环境都产生各自的\omterm{def:b2:magma-emplacement:magma}{岩浆}，因为它们熔化地幔的方式各不相同：
洋脊与热点之下靠减压，俯冲带之上靠加水
（\cref{ch:b2:mantle-melting}）。

\emph{洋中脊}喷发拉斑玄武岩，
即\cref{ch:b2:elements-in-the-earth}中的 MORB，\omterm{def:b2:elements-in-the-earth:partition}{不相容元素}贫乏，
其源区具有亏损的同位素特征。冰岛是位于热点之上的洋脊，
另有自身的拉斑玄武岩，并形成一套直至流纹岩的完整岩系。
% ledger: vp21:series

\emph{洋岛}分阶段生长。夏威夷式火山年轻时喷发碱性熔岩，
随后堆积起由拉斑玄武岩构成的盾状体，估计占其最终体积的 80 至
\qty{95}{\percent}，形成于最初的 1 至 2 百万年内；火山
漂离热点后其熔岩重新变为碱性，并以一长串
稀少的碱性喷发收尾。熔融程度较低时\omterm{def:b2:magma-emplacement:magma}{岩浆}更偏碱性：碱性组分，
作为\omterm{def:b2:elements-in-the-earth:partition}{不相容元素}，在较小的\omterm{def:b2:mantle-melting:fraction}{熔体分数}中稀释较少。同位素方面，\omterm{def:b2:elements-in-the-earth:morb-oib}{洋岛玄武岩}
来自\cref{ch:b2:isotope-geochemistry}中的富集储库。
% ledger: usgs14:hawaii-stages

\emph{火山弧}喷发\omterm{def:b2:magmatic-series:series}{钙碱性系列}。其\omterm{def:b2:elements-in-the-earth:trace}{微量元素}带有
一种独特标记（见下图）：以\omterm{def:b2:elements-in-the-earth:bse}{原始地幔}标准化后，弧玄武岩
富集 Rb、Ba、Th、U、Pb 和 Sr，这些元素由板片加入地幔楔，
相对邻近元素，铌与钽较为贫乏。数据集中的
中位弧玄武岩具有 $\mathrm{Nb}/\mathrm{La} = 0.49$，而 MORB 为 0.97，
\omterm{def:b2:elements-in-the-earth:morb-oib}{洋岛玄武岩}为 1.23。
% ledger: data:b4-basalt-settings

\begin{omfigure}
\centering
\begin{tikzpicture}
\begin{axis}[omrecord, width=0.86\linewidth, height=6.4cm, ymode=log,
  xmin=0.5, xmax=22.5, ymin=1, ymax=300,
  xtick={1,...,22}, xticklabels={Rb,Ba,Th,U,Nb,Ta,La,Ce,Pb,Sr,Nd,Zr,Hf,Sm,Eu,Ti,Gd,Dy,Y,Er,Yb,Lu},
  ylabel={样品 / 原始地幔},
  legend style={font=\scriptsize, at={(0.98,0.98)}, anchor=north east},
  tick label style={font=\scriptsize}, label style={font=\small}]
\addplot[thick, mark=square*, mark size=1.4pt, orange!85!black] table[x=idx, y=oib] {figdata/out/bachelor-2-series-spider.dat};
\addlegendentry{洋岛玄武岩}
\addplot[thick, mark=*, mark size=1.4pt, blue!70!black] table[x=idx, y=morb] {figdata/out/bachelor-2-series-spider.dat};
\addlegendentry{MORB}
\addplot[thick, mark=triangle*, mark size=1.7pt, red!75!black] table[x=idx, y=iab] {figdata/out/bachelor-2-series-spider.dat};
\addlegendentry{岛弧玄武岩}
\end{axis}
\end{tikzpicture}
\par\smallskip{\small 以\omterm{def:b2:elements-in-the-earth:bse}{原始地幔}标准化的中位成分，取自 241 个 MORB、259
个\omterm{def:b2:elements-in-the-earth:morb-oib}{洋岛玄武岩}和 256 个岛弧玄武岩（Vermeesch，2006 年汇编；
\omterm{def:b2:elements-in-the-earth:bse}{原始地幔}取自 McDonough 与 Sun，1995）。弧玄武岩在 Nb
和 Ta 处呈谷值，在 Ba、Pb 和 Sr 处呈峰值；\omterm{def:b2:elements-in-the-earth:morb-oib}{洋岛玄武岩}
则在所有最不相容的元素上均富集。}
\end{omfigure}
% ledger: data:b4-basalt-settings, ms95:pm-trace

在大陆之内，玄武质\omterm{def:b2:magma-emplacement:magma}{岩浆}遇到可被其加热并熔化的厚地壳，
其结果可能是一种\omterm{def:b2:magmatic-series:bimodal}{双峰式火山作用}。

\begin{definition}[双峰式火山作用]\label{def:b2:magmatic-series:bimodal}
\emph{双峰式火山作用}\index{双峰式火山作用}是指在同一岩石省内
喷发两组\omterm{def:b2:magma-emplacement:magma}{岩浆}，即玄武质与流纹质\omterm{def:b2:magma-emplacement:magma}{岩浆}，二者之间
中间成分很少。
\end{definition}

沿黄石热点轨迹的火山作用为双峰式：在东斯内克河平原上，
较老的流纹岩二氧化硅含量约 \qty{75}{\percent}，被
较年轻的玄武岩覆盖，其二氧化硅约 \qty{50}{\percent}。两者之间的空白是
线索：仅凭玄武岩的分异本应留下中间熔岩，而这一空白
指示流纹岩另有来源，如被玄武岩加热的地壳。
% ledger: yvo23:bimodal

\emph{大火成岩省}已见于本科一年级分册，其喷发的玄武岩体量
无可与之比拟的其他构造环境；其系列同样以相同图解判读。

\section{判别图解}\label{sec:b2:magmatic-series:discrimination}

前一节中的规律提示，可直接由分析结果读出构造环境，
岩石学家为此绘制了大量图解。

\begin{definition}[判别图解]\label{def:b2:magmatic-series:discrimination}
一个\emph{判别图解}\index{判别图解}是以化学变量为坐标的图，
变量通常为不易迁移的\omterm{def:b2:elements-in-the-earth:trace}{微量元素}，
图中按已知构造环境岩石的成分划分出若干区域，用以
为成因不明的岩石判定其构造环境。
\end{definition}

经典图解是 Pearce 与 Cann（1973）提出的 Ti--Zr--Y 三角图，针对玄武岩绘制。它
采用蚀变中较少迁移的元素，因此可用于古老的
变质熔岩（其构造环境已不可考）。下图以汇编中三种元素齐全的
538 个玄武岩检验该图。\omterm{def:b2:elements-in-the-earth:morb-oib}{洋岛玄武岩}多落在板内区 D 中，173 个中有 152 个，
\omterm{def:b2:elements-in-the-earth:morb-oib}{洋中脊玄武岩}多落在区 B 中，164 个中有 142 个；但区 B 与弧玄武岩共享，
201 个弧玄武岩中亦有 101 个落于其中。
% ledger: geoplotr:tizry, data:b4-basalt-settings

\begin{omfigure}
\centering
\begin{tikzpicture}
\begin{axis}[width=0.62\linewidth, height=0.55\linewidth, axis equal image,
  hide axis, xmin=-0.04, xmax=1.04, ymin=-0.06, ymax=0.92, clip=false,
  legend style={font=\scriptsize, at={(1.0,0.95)}, anchor=north west}]
\addplot[only marks, mark=*, mark size=0.7pt, blue!70!black, opacity=0.7] table[x=x, y=y] {figdata/out/bachelor-2-series-tizry-morb.dat};
\addlegendentry{MORB}
\addplot[only marks, mark=square*, mark size=0.7pt, orange!85!black, opacity=0.7] table[x=x, y=y] {figdata/out/bachelor-2-series-tizry-oib.dat};
\addlegendentry{洋岛玄武岩}
\addplot[only marks, mark=triangle*, mark size=0.9pt, red!75!black, opacity=0.7] table[x=x, y=y] {figdata/out/bachelor-2-series-tizry-iab.dat};
\addlegendentry{岛弧玄武岩}
\foreach \f in {A,B,C,D}
  \addplot[thick, black, forget plot] table[x=x, y=y] {figdata/out/bachelor-2-series-tizry-\f.dat};
\draw[thick] (axis cs:0,0) -- (axis cs:1,0) -- (axis cs:0.5,0.8660254) -- cycle;
\node[anchor=south] at (axis cs:0.5,0.866) {Ti/100};
\node[anchor=north east] at (axis cs:0,0) {Zr};
\node[anchor=north west] at (axis cs:1,0) {Y $\times$ 3};
\draw[thin] (axis cs:0.33,0.36) -- (axis cs:0.22,0.46) node[anchor=south east, font=\small\bfseries, inner sep=1pt] {D};
\draw[thin] (axis cs:0.36,0.12) -- (axis cs:0.24,0.06) node[anchor=north east, font=\small\bfseries, inner sep=1pt] {C};
\draw[thin] (axis cs:0.56,0.37) -- (axis cs:0.76,0.50) node[anchor=west, font=\small\bfseries, inner sep=1pt] {A};
\draw[thin] (axis cs:0.53,0.21) -- (axis cs:0.72,0.13) node[anchor=north west, font=\small\bfseries, inner sep=1pt] {B};
\end{axis}
\end{tikzpicture}
\par\smallskip{\small Pearce 与 Cann（1973）的 Ti--Zr--Y 图解，附已知构造环境的玄武岩。
A 为岛弧拉斑玄武岩；B 为 MORB、岛弧拉斑玄武岩与
钙碱性玄武岩；C 为钙碱性玄武岩；D 为板内玄武岩。Ti 由
\ce{TiO2} 得出；三种元素均以 ppm 表示。}
\end{omfigure}
% ledger: geoplotr:tizry, data:b4-basalt-settings

\begin{method}[使用判别图解]\label{met:b2:magmatic-series:discrimination}
(1) 检查岩石是否属于该图解所针对的类型，通常为玄武岩，而非
\omterm{def:b2:magmatic-differentiation:cumulate}{堆晶岩}或演化熔岩，其\omterm{def:b2:elements-in-the-earth:trace}{微量元素}反映的结晶作用
不亚于源区。(2) 对于蚀变或变质岩石，宜优先采用蚀变中不活动的元素（Ti、Zr、Y、Nb、重
\omterm{def:b2:elements-in-the-earth:ree}{稀土元素}）。(3) 应投点于一组样品，而非单个样品：
位于边界附近的单个点说明不了问题。(4) 应把结果理解为一种可能性：
各区域是围绕参考样品中的大多数（而非全部）划定的，部分区域
相互重叠。(5) 须以独立证据加以确认：其他图解、同位素、
地质背景。
\end{method}

第 (4) 步可以定量化（见\cref{exo:b2:magmatic-series:8}）。若一玄武岩落在 B 区，且三种环境出现的频率相同，
则它为 MORB 的概率仅约为 0.62。
基于此类汇编数据训练的统计分类器也承袭了同样的局限：玄武岩
保留着其源区与熔融过程的“记忆”，而两种构造环境可能产生相似的
\omterm{def:b2:magma-emplacement:magma}{岩浆}。
% ledger: data:b4-basalt-settings

\begin{inthelab}[检查全岩分析]
在把岩石投到这些图解中的任何一个之前，岩石学家先要检查分析结果。
主量氧化物（例如用 X 射线荧光法测定）之和应接近 100，
再加上挥发分：水和二氧化碳，它们可直接测定，或合并以烧失量
表示，即粉末在强热下失去的质量。烧失量较大
表明岩石已发生蚀变，其碱金属可能已迁移；TAS 图解的方案
会剔除水含量超过 \qty{2}{\percent} 或
二氧化碳含量超过 \qty{0.5}{\percent} 的分析结果。随后将分析结果换算为
不含挥发分、合计为 100 的数值；铁以总 \ce{FeO} 表示，或分为亚铁和高铁，
并将同批次分析的参考物质与其认证值相比较
（见\cref{ch:b2:field-lab-methods-2}）。
\end{inthelab}
% ledger: bgs:tas-use

\section{练习}

\begin{exercise}[$\star$]\label{exo:b2:magmatic-series:1}
在 TAS 图解上命名下列火山岩（\ce{SiO2}、\ce{Na2O}、\ce{K2O}
以 \unit{\percent} 计，已扣除挥发分）：(a) 49.60、2.22、0.51（夏威夷
参考玄武岩 BHVO-2）；(b) 58.5、3.9、1.6；(c) 72.0、4.0、4.2；(d) 47.0、4.5、2.5。
对于 (d)，给出两种可能的岩石名称，并说明什么决定二者之间的选择。
\end{exercise}
% ledger: usgs:bhvo2-major, bgs:tas

\begin{exercise}[$\star$]\label{exo:b2:magmatic-series:2}
某熔岩的分析结果为 \ce{SiO2} 51.60、\ce{TiO2} 1.20、\ce{Al2O3} 16.10、\ce{FeO_T} 9.80、
\ce{MgO} 6.60、\ce{CaO} 8.20、\ce{Na2O} 3.50、\ce{K2O} 1.40、\ce{H2O} 1.40 以及
\ce{CO2} \qty{0.20}{\percent}。(a) 该分析结果是否适用于 TAS 图解？
(b) 将其换算为不含挥发分的数值。(c) 分别给出换算前后
的岩石名称。(d) 它属于低钾、中钾还是高钾？
\end{exercise}
% ledger: bgs:tas-use, bgs:k2o-sio2

\begin{exercise}[$\star$]\label{exo:b2:magmatic-series:3}
在 \omterm{def:b2:magmatic-series:qapf}{QAPF 图解}上命名下列岩石：(a) 粗粒岩石，石英 \qty{25}{\percent}、
碱性长石 \qty{30}{\percent}、斜长石 \qty{35}{\percent}、黑云母 \qty{8}{\percent}，
以及 \qty{2}{\percent} 普通角闪石；(b) 石英 \qty{5}{\percent}、
碱性长石 \qty{5}{\percent}、斜长石 \qty{55}{\percent} 以及
单斜辉石 \qty{35}{\percent}。
\end{exercise}
% ledger: bgs:qapf

\begin{exercise}[$\star$]\label{exo:b2:magmatic-series:4}
借助 \ce{K2O}--\ce{SiO2} 图中的界线，将下列岩石划分为低钾、中钾或高钾：
(a) 50.0、0.40；(b) 56.0、1.50；(c) 62.0、2.90；(d) 66.0、1.00（\ce{SiO2}、
\ce{K2O}、\unit{\percent}）。
\end{exercise}
% ledger: bgs:k2o-sio2

\begin{exercise}[$\star\star$]\label{exo:b2:magmatic-series:5}
(a) 对一玄武岩计算 $A$、$F$ 和 $M$，其 \ce{FeO_T} 为 14.5、\ce{MgO} 为 4.5 和
\ce{Na2O + K2O} 为 \qty{3.2}{\percent}；对一安山岩，其数值为 6.1、3.8 和 5.05。
(b) 划分线经过 $(F, A, M) = (54.5, 23.0, 22.5)$、$(52.5, 30.0,
17.5)$ 和 $(49.6, 35.0, 15.4)$，其中第一点为其最富铁的点；拉斑一侧
即 F 所在一侧。每种岩石落在哪一侧？(c) 哪种
矿物在早期结晶，使安山岩留在原处？
\end{exercise}
% ledger: geoplotr:afm, vp21:series

\begin{exercise}[$\star\star$]\label{exo:b2:magmatic-series:6}
汇编数据中 MORB 的中位数含 Nb \qty{4}{ppm}、La \qty{4.12}{ppm}，
\omterm{def:b2:elements-in-the-earth:morb-oib}{洋岛玄武岩}的中位数为 38 和 30.8，岛弧玄武岩的中位数为 3 和 6.1。取
\omterm{def:b2:elements-in-the-earth:bse}{原始地幔}值 Nb 0.658 和 La \qty{0.648}{ppm}，计算各自的标准化值
以及比值 $\mathrm{Nb_N}/\mathrm{La_N}$。哪一种具有铌异常？原因何在？
\end{exercise}
% ledger: data:b4-basalt-settings, ms95:pm-trace

\begin{exercise}[$\star\star$]\label{exo:b2:magmatic-series:7}
Ti 质量为 \ce{TiO2} 质量的 $0.5993$。计算三个玄武岩的 Ti--Zr--Y 图坐标
（Ti/100、Zr、$3\mathrm Y$，重新换算为 100）：(a) \ce{TiO2}
\qty{1.6}{\percent}、Zr 110、Y \qty{32}{ppm}；(b) 2.8、220、29；(c) 0.8、50、20。在图中标出
它们的位置，并说明每个位置能使你得出什么结论。
\end{exercise}
% ledger: geoplotr:tizry

\begin{exercise}[$\star\star$]\label{exo:b2:magmatic-series:8}
在本章的 Ti--Zr--Y 检验中，B 区中含有 164 个 MORB 中的 142 个、173 个\omterm{def:b2:elements-in-the-earth:morb-oib}{洋岛玄武岩}中的 4 个
以及 201 个岛弧玄武岩中的 101 个。(a) 各组落入 B 区的比例各为多少？
(b) 若三种构造环境出现频率相同，则投点于 B 区的
玄武岩为 MORB 的概率为多少？(c) 仅根据 B 区，能得出什么结论？
\end{exercise}
% ledger: data:b4-basalt-settings

\begin{exercise}[$\star\star$]\label{exo:b2:magmatic-series:9}
一流纹岩的 \ce{Na2O} 为 5.5、\ce{K2O} 为 4.5、\ce{Al2O3} 为 \qty{10.5}{\percent}（$M$ =
61.98、94.20 和 \qty{101.96}{g/mol}）。计算其分子比 $(\ce{Na2O +
K2O})/\ce{Al2O3}$。(a) 它是否为过碱性？(b) 碱性长石每个 \ce{Al2O3}
只容纳一个碱金属氧化物，则有多大比例的碱金属无法进入长石？
\end{exercise}
% ledger: bgs:tas-use

\begin{exercise}[$\star\star\star$]\label{exo:b2:magmatic-series:10}
在对数比命题的模型中，取 $\lambda_A = 0.4$ 与 $\lambda_M =
0.8$。(a) 分别计算 $\lambda_F = 0.5$ 与 $\lambda_F =
0.7$ 时趋势线的斜率。(b) 冰岛岩石的斜率为 $-1.10$，喀斯喀特山脉岩石的斜率为
$-0.72$，各自隐含的 $\lambda_F$ 为多少？(c) 哪一系列失铁更快，
并由何种矿物加以解释？(d) 沿各趋势线，$F/M$ 如何变化？
\end{exercise}
% ledger: vp21:model

\begin{exercise}[$\star\star\star$]\label{exo:b2:magmatic-series:11}
一种简化的标准矿物计算法。以氧化物的摩尔数计。将全部 \ce{K2O} 归入正长石
（\ce{K2O.Al2O3.6SiO2}），将全部 \ce{Na2O} 归入钠长石（\ce{Na2O.Al2O3.6SiO2}），
剩余的 \ce{Al2O3} 与尽可能多的 \ce{CaO} 归入钙长石（\ce{CaO.Al2O3.2SiO2}），
剩余的 \ce{CaO} 与尽可能多的 $(\ce{Mg,Fe})\ce{O}$ 归入透辉石
（\ce{CaO.(Mg,Fe)O.2SiO2}），其余的 $(\ce{Mg,Fe})\ce{O}$ 归入紫苏辉石
（\ce{(Mg,Fe)O.SiO2}）。若二氧化硅不足，则将紫苏辉石转为橄榄石
（\ce{2(Mg,Fe)O.SiO2}）；若仍不足，则将钠长石转为霞石
（\ce{Na2O.Al2O3.2SiO2}）。将其应用于 (a) BHVO-2：\ce{SiO2} 49.60，\ce{Al2O3} 13.44，
\ce{FeO_T} 11.15，\ce{MgO} 7.257，\ce{CaO} 11.40，\ce{Na2O} 2.219，\ce{K2O}
\qty{0.513}{\percent}；(b) 一种玄武岩，其成分分别为 45.5、15.0、11.5、8.5、10.0、3.8 和 1.5（摩尔
质量依次为 \ce{SiO2} 60.08，\ce{Al2O3} 101.96，\ce{FeO} 71.84，\ce{MgO} 40.30，\ce{CaO}
56.08，\ce{Na2O} 61.98，\ce{K2O} \qty{94.20}{g/mol}）。两者是否各自\omterm{def:b2:magmatic-series:saturation}{硅饱和}？
各玄武岩分别属于哪一系列？
\end{exercise}
% ledger: usgs:bhvo2-major, bgs:series-terms

\begin{exercise}[$\star\star\star$]\label{exo:b2:magmatic-series:12}
在汇编数据中，$\mathrm{Nb}/\mathrm{La} < 0.7$ 对 141 个弧玄武岩中的 116 个成立，
对 101 个 MORB 中的 21 个成立，对 137 个\omterm{def:b2:elements-in-the-earth:morb-oib}{洋岛玄武岩}中的 3 个成立（两者均已分析）。(a) 计算
这三个比例。(b) 若三种构造环境比例相等，则
$\mathrm{Nb}/\mathrm{La} < 0.7$ 的玄武岩为弧玄武岩的概率是多少？(c)
对某地区重做 (b)，其玄武岩中 \qty{60}{\percent} 为 MORB，\qty{30}{\percent}
为\omterm{def:b2:elements-in-the-earth:morb-oib}{洋岛玄武岩}，\qty{10}{\percent} 为弧玄武岩。(d) (c) 的结果对\omterm{def:b2:magmatic-series:discrimination}{判别图解}
的使用有何启示？
\end{exercise}
% ledger: data:b4-basalt-settings

\section{问题：一座火山，三种熔岩}

\begin{problem}[{周末问题 --- 由三份分析确定一座火山的构造环境及
其遗留的晶体}]
\label{pb:b2:magmatic-series:1}
某（虚构的）层状火山的三种熔岩，从最古老的侧翼流动熔岩到最年轻的熔丘，
其成分（\unit{\percent}）为：

\smallskip\noindent{\small
\begin{tabular}{@{}lrrrrrrrrrr@{}}
 & \ce{SiO2} & \ce{TiO2} & \ce{Al2O3} & \ce{FeO_T} & \ce{MgO} & \ce{CaO} & \ce{Na2O} & \ce{K2O} & \ce{H2O} & \ce{CO2}\\
熔岩 1 & 50.4 & 0.95 & 17.6 & 8.6 & 7.9 & 10.2 & 2.9 & 0.62 & 0.70 & 0.10\\
熔岩 2 & 58.6 & 0.80 & 17.4 & 6.1 & 3.8 & 6.9 & 3.9 & 1.15 & 1.10 & 0.05\\
熔岩 3 & 65.3 & 0.55 & 16.2 & 4.0 & 1.6 & 4.3 & 4.4 & 1.85 & 1.20 & 0.05\\
\end{tabular}}

\smallskip\noindent 另以 ppm 计：Rb 10、20、32；Th 1.1、2.1、3.4；Nb 2.0、3.6、5.6；
La 7、13、20；Sr 520、480、380；Zr 70、125、190；Y 18、22、24。\omterm{def:b2:elements-in-the-earth:bse}{原始地幔}：Nb
0.658，La 0.648，Th \qty{0.0795}{ppm}。

\noindent\textbf{第一部分 --- 熔岩命名。}
\begin{enumerate}
\item 将每份分析的各项求和。三份分析是否都可用于 TAS 图解？
\item 将熔岩 1 换算为无挥发分的成分，给出换算因子、\ce{SiO2} 与
\ce{Na2O + K2O}。
\item 对熔岩 2 和 3 做同样的计算。
\item 在 TAS 图解上为三种熔岩命名。
\item 若不作换算，名称是否会改变？
\item 一块成分与熔岩 2 相同的粗粒岩石，结晶于该火山的
根部。哪一种分类方案可为其命名？你预期它落在哪个区域？
\end{enumerate}

\noindent\textbf{第二部分 --- 属于哪一系列？}
\begin{enumerate}[resume]
\item 将每种熔岩划分为低钾、中钾或高钾类型。
\item 计算各熔岩（无挥发分）的 $A$、$F$、$M$。
\item 分界线经过 $(F, A, M) = (42.8, 17.1, 40.1)$、$(46.9, 18.1,
35.0)$、$(52.5, 30.0, 17.5)$、$(49.6, 35.0, 15.4)$、$(40.8, 50.0, 9.2)$ 和 $(37.9,
55.0, 7.1)$，拉斑玄武岩区位于 F 一侧。在相邻两点之间，将 $F$ 对 $A$ 线性插值，
各熔岩落在分界线的哪一侧？
\item 计算各熔岩的 \ce{FeO_T}/\ce{MgO}。是否存在铁的富集？
\item 该火山属于哪一系列？可由哪种矿物和何种条件
加以解释？
\end{enumerate}

\noindent\textbf{第三部分 --- 属于哪种构造环境？}
\begin{enumerate}[resume]
\item 将熔岩 1 的 Nb、La 和 Th 相对\omterm{def:b2:elements-in-the-earth:bse}{原始地幔}作标准化。
\item 计算熔岩 1 的 $\mathrm{Nb_N}/\mathrm{La_N}$，并与
\cref{exo:b2:magmatic-series:6} 中的中位数作比较。
\item 计算 $\mathrm{Th_N}/\mathrm{Nb_N}$。是什么携带了钍？
\item 计算熔岩 1 的 Ti--Zr--Y 坐标（\ce{TiO2} 取无挥发分分析的
值）。其所在区域说明了什么？为何不对熔岩 2 和 3 作图？
\item 作出结论：这座火山形成于何种构造环境？
\end{enumerate}

\noindent\textbf{第四部分 --- 结晶了多少？}
\begin{enumerate}[resume]
\item 钾不进入任何早期矿物。由熔岩 1 与 2 的无挥发分 \ce{K2O}，
计算熔岩 2 形成时残留液体的分数 $F$，若其
由熔岩 1 经\omterm{def:b2:magmatic-differentiation:crystallisation}{分离结晶}形成。
\item 对铷和钍做同样的计算。
\item 锆的整体分配系数 $D \approx 0.1$，利用瑞利定律
由它估计 $F$。
\item 锶从熔岩 2 降至熔岩 3。由熔岩 2 与 3 之间的钾计算 $F$，
并据此估计该区间锶的整体分配系数 $D$。哪种矿物
应对此负责？
\item 从熔岩 1 到熔岩 3，有多大比例的液体结晶了？
\item 若\omterm{def:b2:magma-emplacement:magma}{岩浆}同化了富钾的花岗质地壳，钾的估计值会
偏高还是偏低？
\item 对第四部分中熔岩 1 至 2 的四个估计值取平均，并给出
每 \qty{100}{kg} 玄武岩为形成安山岩而遗留的晶体质量，
即本题的关键数值。
\end{enumerate}
\end{problem}
% ledger: bgs:tas-use, bgs:tas, bgs:k2o-sio2, geoplotr:afm, geoplotr:tizry, ms95:pm-trace
